\subsection{单环}

命题 1. —— 设 A 是一个非零环。下列条件等价：

\begin{enumerate}
\def\labelenumi{(\roman{enumi})}
\item
  A-模 \(A_{s}\) 是同型的。
\item
  环 A 是左阿廷的，并且 A 的每个双边理想都等于 0 或 A。
\item
  环 A 是左阿廷的，并且存在一个既单又忠实的左 A-模 S。
\end{enumerate}

若这些条件满足，则 A-模 \(A_{s}\) 具有有限长度且是 S 型同型的，并且每个单 A-模都同构于 S。

我们来证 (i) 蕴含 (ii)。设 (i) 满足。那么有限生成 A-模 \(A_{s}\) 是半单的，因而具有有限长度且是阿廷的（VIII, 第 71 页，命题 10）；于是环 A 是左阿廷的。左 A-模 \(A_{s}\) 的自同态是由 A 的元素作的右乘。由于左 A-模 \(A_{s}\) 是同型的，由 VIII, 第 86 页的命题 6, b)，(A, A)-双模 \(_{s}A_{d}\) 是单的。\(_{s}A_{d}\) 的子双模就是 A 的双边理想，故 (i) 蕴含 (ii)。

我们来证 (ii) 蕴含 (iii)。环 A 不约化为 0；于是存在一个单 A-模 S。S 的零化理想是 A 的一个真双边理想。若设 (ii) 满足，则该零化理想等于 0。于是 A-模 S 是忠实的，故 (ii) 蕴含 (iii)。

我们来证 (iii) 蕴含 (i)。设 (iii) 满足。那么存在一个整数 \(m \geqslant 1\)，使 \(A_{s}\) 同构于 \(S^{m}\) 的一个子模（VIII, 第 50 页，命题 5, a)）。由于 \(S^{m}\) 是 S 型的同型 A-模，\(A_{s}\) 亦然（VIII, 第 61 页，命题 2）；故 (iii) 蕴含 (i)。

设条件 (i)–(iii) 满足。我们上面已看到 A-模 \(A_{s}\) 具有有限长度且是 S 型同型的。每个单左 A-模都同构于 \(A_{s}\) 的一个商，因而同构于 S。

定义 1. —— 若环 A 满足命题 1 的等价条件 (i)、(ii) 与 (iii)，则称它为单环。设 K 是一个交换域（域）；若一个 K-代数的底环是单环，则称该 K-代数为单的。

注. —— 1) 回忆由 II, §9, No.~1, 第 115 页的定理 1，下列性质等价：

\begin{enumerate}
\def\labelenumi{(\roman{enumi})}
\item
  A-模 \(A_{s}\) 是单的。
\item
  环 A 不约化为 0，并且 A 的左理想只有 0 与 A。
\item
  环 A 是体。于是，由命题 1（条件 (ii)），交换单环无非就是交换域（域）。
\end{enumerate}

\pandocbounded{\includegraphics[keepaspectratio]{images/part001_205e3f38.jpg}}

\begin{enumerate}
\def\labelenumi{\arabic{enumi})}
\setcounter{enumi}{1}
\tightlist
\item
  若环 A 不约化为 0 且其仅有的双边理想是 0 与 A，我们有时称 A 是拟单环。若 A 具有一个忠实单模，则称 A 是本原的。由命题 1，每个单环都是拟单环。由于每个非零环都具有一个单模，且单模的零化理想是双边理想，可见每个拟单环都是本原的。然而，存在不是单环的拟单环以及不是拟单环的本原环（VIII, 第 128 页，习题 2）；这样的环不是左阿廷的。
\end{enumerate}

定理 1（韦德伯恩）. —— 一个环是单环当且仅当它同构于一个矩阵环 \(\mathbf{M}_r(\mathrm{D})\)，其中 \(r \geqslant 1\) 是整数，D 是一个体。

引理 1. —— 设 A 是一个单环，S 是一个单左 A-模，D 是体 \(\operatorname{End}_{\mathrm{A}}(\mathrm{S})\) 的反环。那么 S 是一个可逆 (A,D)-双模。它也是 D 上的一个有限维右向量空间，并且映射 \(a \mapsto a_{S}\) 是从 A 到 \(\operatorname{End}_{\mathrm{D}}(\mathrm{S})\) 的环同构。

由命题 1，A-模 \(A_{s}\) 具有有限长度且是 S 型同型的。于是存在一个整数 \(m \geqslant 1\)，使 A-模 \(A_{s}\) 与 \(S^{m}\) 同构。那么 A-模 S 是投射的且有限生成的。它是生成模（VIII, 第 80 页，定理 1），而引理 1 由 VIII, 第 101 页的定理 1 的 (ii)⇒(i) 与 (ii)⇒(iii) 应用于 \((\mathrm{A},\mathrm{D})_{\mathbf{Z}}\)-双模 S 得出。

引理 2. —— 设 D 是一个体，V 是 D 上一个维数有限 \(r \geqslant 1\) 的右向量空间。那么 V 是环 \(\mathrm{E} = \mathrm{End}_{\mathrm{D}}(\mathrm{V})\) 上的单模，且它的交换子等于 \(D_{V}\)。环 E 是单环，其左长度等于 r。

我们知道 V 是一个单 E-模（VIII, 第 45 页，例 3），且其交换子等于 \(D_{V}\)（VIII, 第 82 页，推论 1）。设 \((x_{i})_{1\leqslant i\leqslant r}\) 是 V 在体 D 上的一个基。映射 \(u\mapsto(u(x_{i}))_{1\leqslant i\leqslant r}\) 是从 E-模 \(E_{s}\) 到 E-模 \(V^{r}\) 的同构；于是 E-模 \(E_{s}\) 是长度为 r 的同型模，故环 E 是单环。

现在证明定理 1。回忆（II, §10, No.~7, 第 349 页）环 \(\mathbf{M}_r(\mathrm{D})\) 可以等同于右 D-向量空间 \(\mathrm{D}_d^r\) 的自同态环；此外，体 D 上每个维数为 \(r\) 的右向量空间都同构于 \(\mathrm{D}_d^r\)（II, §7, No.~1, 第 292 页）。于是定理 1 由引理 1 与引理 2 得出。

注 3. —— 设 A 是一个单环，S 是一个单 A-模，D 是体 \(\operatorname{End}_{\mathrm{A}}(\mathrm{S})\) 的反环。那么 A-模 \(A_{s}\) 具有有限长度，且 \(\dim_{\mathrm{D}}(\mathrm{S})\) 等于 \(\operatorname{long}(\mathrm{A})\)。事实上，由引理 1，环 A 同构于 \(\operatorname{End}_{\mathrm{D}}(\mathrm{S})\)；然后应用引理 2 即可。

推论 1. —— a) 单环的中心是一个体。

\begin{enumerate}
\def\labelenumi{\alph{enumi})}
\setcounter{enumi}{1}
\item
  单环的反环是单环。
\item
  单环的左长度等于其右长度。
\end{enumerate}

设 D 是一个体，Z 是它的中心，V 是 D 上一个维数有限 \(r \geqslant 1\) 的右向量空间。我们用 E 表示 V 的自同态环。

由 VIII, 第 83 页的推论 2，映射 \(z \mapsto z_{V}\) 是从 Z 到 E 的中心的同构。断言 a) 得证。V 的对偶 \(V^{*}\) 是 D 的反体 \(D^{o}\) 上的右向量空间，其维数等于 r。映射 \(u \mapsto {}^{t}u\) 是从 E 的反环 \(E^{o}\) 到环 \(\operatorname{End}_{\mathrm{D}^{\circ}}(\mathrm{V}^{*})\) 的同构。于是环 \(E^{o}\) 是单环，且环 E 与 \(E^{o}\) 有相同的左长度，都等于 r（引理 2）。

推论 2. —— 设 r 与 \(r'\) 是严格正的整数，D 与 \(D'\) 是体。环 \(\mathbf{M}_{r}(\mathrm{D})\) 与 \(\mathbf{M}_{r'}(\mathrm{D}')\) 同构当且仅当我们有 \(r = r'\) 且体 D 与 \(D'\) 同构。

该条件显然是充分的。

反之，设环 \(\mathrm{B} = \mathbf{M}_{r}(\mathrm{D})\) 与 \(\mathrm{B}' = \mathbf{M}_{r'}(\mathrm{D}')\) 同构。由于 r 是 \(B_{s}\) 的长度而 \(r'\) 是 \(B'_{s}\) 的长度（引理 2），我们有 \(r = r'\)。此外，B 与 D 森田等价，\(B'\) 与 \(D'\) 森田等价（VIII, 第 102 页，

例 1）。于是体 D 与 D' 是森田等价的，因而同构（VIII, 第 111 页，命题 13, c)）。

推论 3. —— 设 K 是一个交换域（域），A 是一个底环为单环的有限次数 K-代数。存在一个整数 r 与一个 K 上有限次数的 K-代数 D，D 是体，使 A 同构于 \(\mathrm{M}_{r}(\mathrm{D})\)。特别地，若 K 是代数闭的，则 A 同构于 K 上的一个矩阵代数。

设 S 是一个单左 A-模；它是 K 上的一个有限维向量空间。于是它的交换子是 K 上有限次数的代数。从而第一条断言由引理 1 得出。若 K 是代数闭的，则由 VIII, 第 47 页的定理 1 有 D = K。

注 4. —— 设 K 是一个代数闭的交换域（域），A 是 K 上一个有限次数的代数。代数 A 是单的当且仅当存在一个整数 \(n \geqslant 1\)，使 A 同构于 \(\mathrm{M}_{n}(\mathrm{K})\)。此时它的中心同构于 K。

